\documentclass[11pt,a4paper]{article}

\usepackage[T1]{fontenc}
\usepackage[utf8]{inputenc}
\usepackage{lmodern}
\usepackage[margin=1.05in]{geometry}
\usepackage{amsmath,amssymb}
\usepackage{booktabs}
\usepackage{xcolor}
\usepackage{enumitem}
\usepackage{listings}
\usepackage{titlesec}
\usepackage[hidelinks]{hyperref}

\definecolor{accent}{HTML}{17547E}
\definecolor{signal}{HTML}{A93C1C}
\definecolor{good}{HTML}{2A6647}
\definecolor{muted}{HTML}{5F6A78}
\definecolor{rule}{HTML}{C2CCD8}
\definecolor{panel}{HTML}{F1F4F8}

\titleformat{\section}{\large\bfseries}{\color{accent}\thesection}{0.7em}{}
\titleformat{\subsection}{\normalsize\bfseries}{\color{accent}\thesubsection}{0.6em}{}

\lstset{
  basicstyle=\ttfamily\scriptsize,
  commentstyle=\color{muted},
  breaklines=true,
  columns=fullflexible,
  keepspaces=true,
  frame=single,
  rulecolor=\color{rule},
  backgroundcolor=\color{panel},
  framesep=5pt,
  xleftmargin=0pt,
  aboveskip=1em,
  belowskip=1em,
  showstringspaces=false,
}

\newcommand{\defect}[2]{%
  \par\medskip\noindent
  {\color{signal}\rule{\linewidth}{1.2pt}}\par\vspace{-0.3em}
  \noindent{\color{signal}\bfseries\sffamily #1}\quad{\bfseries #2}\par
  \vspace{0.2em}\noindent{\color{rule}\rule{\linewidth}{0.4pt}}\par\vspace{0.4em}
}
\newcommand{\defectend}{\par\vspace{0.2em}\noindent{\color{signal}\rule{\linewidth}{1.2pt}}\par\medskip}

\newcommand{\ok}{{\color{good}\textsf{\footnotesize verified}}}
\newcommand{\fixed}{{\color{good}\textsf{\footnotesize fixed}}}
\newcommand{\bad}[1]{{\color{signal}\textsf{\footnotesize #1}}}
\newcommand{\warn}[1]{{\color{accent}\textsf{\footnotesize #1}}}

\newcommand{\eqp}{\mathrm{eq}}
\newcommand{\Fold}{\mathrm{Fold}}
\newcommand{\RS}{\mathrm{RS}}
\newcommand{\CRS}{\mathrm{CRS}}
\newcommand{\pow}{\mathrm{pow}}
\newcommand{\err}{\mathrm{err}}

\setlength{\emergencystretch}{2.5em}
\setlist[itemize]{itemsep=0.25em,topsep=0.35em,leftmargin=1.4em}
\setlist[enumerate]{itemsep=0.25em,topsep=0.35em,leftmargin=1.8em}

\title{\vspace{-2em}\textbf{WHIR on Bitcoin Script}\\[0.35em]
  {\large\normalfont A formal account of the STIR and WHIR proximity tests,\\
  and a review of \texttt{bitcoin-stark-verifier} against them}}
\author{}
\date{\small Subject: \texttt{AppliedPQC/bitcoin-stark-verifier}
  \quad$\cdot$\quad Sources: ePrint 2024/390, 2024/1586
  \quad$\cdot$\quad 98 tests passing
  \quad$\cdot$\quad D1--D3 fixed; queries composed, threaded and closed}

\begin{document}
\maketitle
\thispagestyle{empty}

\begin{center}
\begin{minipage}{0.94\linewidth}
\small
\noindent{\color{accent}\rule{\linewidth}{1pt}}\par\vspace{0.4em}
\noindent\textbf{\sffamily\footnotesize SUMMARY OF FINDINGS}\par\vspace{0.5em}
\begin{itemize}[leftmargin=1.2em]
  \item \textbf{The primitives are sound and genuinely verified.} Poseidon2, the duplex sponge, the
  challenger, the sumcheck round and multilinear evaluation are each checked against a Rust
  transcription of Plonky3, with negative cases. A real Plonky3 WHIR proof's Merkle opening
  authenticates in script, and a tampered leaf is rejected.
  \item \textbf{Three defects sat above them}, all in the composition rather than the arithmetic:
  Fiat--Shamir challenges were never derived (\textbf{D1}), the closing identity was a tautology
  (\textbf{D2}), and the composed \texttt{verify()} was a cost model that was never executed
  (\textbf{D3}). All three are fixed, with adversarial tests.
  \item \textbf{The cost is hashing, and almost nothing else.} The example configuration is 2162
  permutations, of which \textbf{97.1\%} is query work --- Merkle paths and the leaf hash that binds
  a row to them --- against 2.9\% for the whole transcript.
  \item \textbf{The verifier ends on its identity, and derives everything in it.} The queries are
  composed, their roots bound to the absorbed commitment, the folding randomness accumulated, and
  the closing check $\texttt{claimed} = w(\boldsymbol R)\cdot\hat f_{M}(\boldsymbol
  r_{\mathrm{fin}})$ emitted with its constraints \emph{derived from the transcript} rather than
  supplied --- for 7.5\% of the verifier and no permutation at all. Only the statement is supplied,
  because a statement is what a claim \emph{is}.
  \item \textbf{Two configuration levers remain unused}, worth roughly 20\% of script weight
  between them: \texttt{pow\_bits = 0} where both reference implementations use $\approx$22, and
  the choice of soundness regime, currently the conjectural one.
\end{itemize}
\vspace{0.3em}\noindent{\color{accent}\rule{\linewidth}{1pt}}
\end{minipage}
\end{center}

\vspace{1em}

\begin{center}{\large\bfseries Part I\quad The algorithm, formally}\end{center}
\vspace{-0.5em}
\noindent{\color{rule}\rule{\linewidth}{0.6pt}}

\section{Setting and notation}

Fix a finite field $\mathbb{F}$ with $\mathrm{char}(\mathbb{F})\neq 2$ and a \emph{smooth}
evaluation domain $\mathcal{L}\subseteq\mathbb{F}^{*}$ --- a multiplicative coset whose order is a
power of two. For a degree bound $d$,
\[
  \RS[\mathbb{F},\mathcal{L},d]
  = \bigl\{\, f:\mathcal{L}\to\mathbb{F} \;\bigm|\; \exists\,\hat g\in\mathbb{F}^{<d}[X],\;
    f(x)=\hat g(x)\ \forall x\in\mathcal{L} \,\bigr\},
  \qquad \rho := d/|\mathcal{L}| .
\]
When $d=2^{m}$ the code has an equivalent multilinear reading, and this is the view both protocols
use. Writing $\pow(x,m)=(x^{2^{0}},\dots,x^{2^{m-1}})$,
\[
  \RS[\mathbb{F},\mathcal{L},m]
  = \bigl\{\, f \;\bigm|\; \exists\,\hat f\in\mathbb{F}^{<2}[X_{1},\dots,X_{m}],\;
    f(x)=\hat f(\pow(x,m)) \,\bigr\}.
\]
$\Delta(f,g)$ is fractional Hamming distance and $\Lambda(\mathcal{C},f,\delta)$ the list. The
equality polynomial is
\[
  \eqp(\boldsymbol X,\boldsymbol Y)=\prod_{i=1}^{m}\bigl(X_{i}Y_{i}+(1-X_{i})(1-Y_{i})\bigr),
  \qquad
  \hat f(\boldsymbol z)=\sum_{\boldsymbol b\in\{0,1\}^{m}}\hat f(\boldsymbol b)\cdot
  \eqp(\boldsymbol b,\boldsymbol z).
\]

\section{Constrained Reed--Solomon codes}\label{sec:crs}

WHIR's object of study is not $\RS$ but a subcode carrying one sumcheck-shaped constraint:
\[
  \CRS[\mathbb{F},\mathcal{L},m,\hat w,\sigma]
  = \Bigl\{\, f\in\RS[\mathbb{F},\mathcal{L},m] \;:\;
  \textstyle\sum_{\boldsymbol b\in\{0,1\}^{m}}\hat w(\hat f(\boldsymbol b),\boldsymbol b)=\sigma
  \,\Bigr\}.
  \tag{Def.\ 1}
\]
\subsection*{The weight polynomial}

Definition 1 leans on $\hat w$ without saying what it is, and it is worth being
precise, because every later step is a statement about $\hat w$ and almost all of the
implementation work is bookkeeping for it.

$\hat w$ is a polynomial in $1+m$ variables,
\[
  \hat w(Z,\boldsymbol X)=\hat w(Z,X_{1},\dots,X_{m}),
\]
whose first argument $Z$ is a \emph{placeholder for a value of the polynomial being
tested} and whose remaining arguments $\boldsymbol X$ range over positions on the
hypercube. Substituting $Z\mapsto\hat f(\boldsymbol b)$ and $\boldsymbol
X\mapsto\boldsymbol b$ and summing over $\boldsymbol b\in\{0,1\}^{m}$ turns it into a
single field element, and Definition 1 asks that element to be $\sigma$. So $\hat w$ is
not a constraint on $f$ pointwise --- it is the \emph{shape} of one aggregate linear
statement about $f$, and $\sigma$ is its claimed value.

Two restrictions make it usable rather than merely general.

\smallskip
\noindent\textbf{(i) Linear in $Z$.} Every $\hat w$ WHIR forms has the shape
\[
  \hat w(Z,\boldsymbol X)=Z\cdot w(\boldsymbol X),
  \qquad w \text{ multilinear in } \boldsymbol X,
\]
so the sum $\sum_{\boldsymbol b}\hat f(\boldsymbol b)\,w(\boldsymbol b)$ is a sumcheck
over a \emph{product of two multilinears}. That is what fixes the degree of every round
polynomial at two, which in turn is why a round sends three values --- and why, with
$\hat h(0)+\hat h(1)=\sigma$ pinning one of them, only $\hat h(0)$ and $\hat h(\infty)$
travel. The entire cost model of \S\ref{sec:cost} follows from this one line.

\smallskip
\noindent\textbf{(ii) Closed under the recursion.} \S6's update
\[
  \hat w'(Z,\boldsymbol X)=\hat w(Z,\boldsymbol\alpha,\boldsymbol X)
  + Z\cdot\sum_{i=0}^{t}\gamma^{\,i+1}\eqp(\boldsymbol z_{i},\boldsymbol X)
\]
is linear in $Z$ whenever $\hat w$ is, since binding $\boldsymbol\alpha$ does not touch
$Z$ and the added term is visibly $Z$ times something. So $\hat w'$ is again of the form
$Z\cdot w'(\boldsymbol X)$, and the protocol never leaves the class it started in.

\smallskip
Unfolding (ii) across all $M$ rounds gives the concrete object the verifier ends up
holding:
\[
  w(\boldsymbol X)=\sum_{j} c_{j}\cdot\eqp(\boldsymbol z_{j},\boldsymbol X),
\]
a \textbf{weighted sum of equality polynomials} --- that is, nothing more than a list of
pairs $(c_{j},\boldsymbol z_{j})$. Every constraint the protocol ever accumulates, from
the initial evaluation query through each round's out-of-domain answer and $t$ shift
answers, is one entry in that list, and the coefficients $c_{j}$ are powers of the
batching challenges. \S\ref{sec:constrain} is about carrying that list, and \S\ref{sec:compose} is about what it
costs to evaluate it.

\subsection*{Evaluation as a special case}

The point of the definition is that a polynomial \emph{evaluation query} is a special case. Taking
$\hat w(Z,\boldsymbol X)=Z\cdot\eqp(\boldsymbol X,\boldsymbol z)$ gives
$\sum_{\boldsymbol b}\hat w(\hat f(\boldsymbol b),\boldsymbol b)=\hat f(\boldsymbol z)$, so
membership with that weight \emph{is} the claim $\hat f(\boldsymbol z)=\sigma$. A proximity test for
$\CRS$ codes is therefore simultaneously a low-degree test and a multilinear polynomial commitment
--- which is why WHIR replaces FRI, STIR and BaseFold at once rather than sitting beside them.

\section{Folding}

For $\alpha\in\mathbb{F}$ and $y=x^{2}$,
\[
  \Fold(f,\alpha)(x^{2}) \;=\; \frac{f(x)+f(-x)}{2}
  \;+\; \alpha\cdot\frac{f(x)-f(-x)}{2x},
\]
extended to $\boldsymbol\alpha=(\alpha_{1},\dots,\alpha_{k})$ by iteration. The structural fact is
that folding is multilinear partial evaluation seen through the code: if
$f\in\RS[\mathbb{F},\mathcal{L},m]$ then
\[
  \Fold(f,\boldsymbol\alpha)\in\RS[\mathbb{F},\mathcal{L}^{(2^{k})},m-k],
  \qquad
  \hat g(X_{k+1},\dots,X_{m})=\hat f(\boldsymbol\alpha,X_{k+1},\dots,X_{m}).
  \tag{Claim 4.15}
\]
Three quantities move together, and it is worth reading $m-k$ as the count it is: binding the first
$k$ variables to $\boldsymbol\alpha$ leaves $X_{k+1},\dots,X_{m}$ free, so the arity falls from $m$
to $m-k$, the degree from $2^{m}$ to $2^{m-k}$, and the domain from $\mathcal{L}$ to
$\mathcal{L}^{(2^{k})}$ of size $|\mathcal{L}|/2^{k}$.

That third movement is the one to watch, because of what it does to the rate:
\[
  \rho=\frac{2^{m}}{|\mathcal{L}|}
  \qquad\longrightarrow\qquad
  \rho'=\frac{2^{m-k}}{|\mathcal{L}|/2^{k}}=\frac{2^{m}}{|\mathcal{L}|}=\rho .
\]
\textbf{Folding alone does not change the rate.} Claim 4.15 is the structural half of the story and
says nothing about cost. The rate reduction of \S4 comes from somewhere else entirely --- from the
prover re-committing $g$ on a domain that is merely \emph{halved} rather than divided by $2^{k}$ ---
which is why that section is about a choice and this one is about an identity.

\section{Rate is the design variable}

This is the whole idea behind both protocols, and it is a single line of arithmetic.

\begin{center}\small
\begin{tabular}{@{}lllll@{}}
\toprule
\textbf{Protocol} & \textbf{Degree} & \textbf{Domain} & \textbf{Rate} & $\log_{2}(1/\rho)$ per round\\
\midrule
FRI  & $d\to d/k$ & $|\mathcal{L}|\to|\mathcal{L}|/k$ & $\rho'=\rho$ & unchanged\\
STIR & $d\to d/k$ & $|\mathcal{L}|\to|\mathcal{L}|/2$ & $\rho'=(2/k)\rho$ & $+\log_{2}k-1$\\
WHIR & $m\to m-k$ & $|\mathcal{L}|\to|\mathcal{L}|/2$ & $\rho'=2^{1-k}\rho$ & $+(k-1)$\\
\bottomrule
\end{tabular}
\end{center}

\medskip
\noindent\fbox{\begin{minipage}{0.96\linewidth}\small
\textbf{A notational trap.} STIR's $k$ is the degree \emph{divisor} (a power of two, $k\ge4$);
WHIR's $k$ is the \emph{number of variables} folded. The two agree under
$k_{\mathrm{STIR}}=2^{\,k_{\mathrm{WHIR}}}$. The recurrence that appears in code is WHIR's.
\end{minipage}}
\medskip

The number of queries needed in round $i$ to reach round-by-round error $2^{-\lambda}$ is
$t_{i}=\lceil \lambda / -\log_{2}(\rho_{i}+\eta_{i})\rceil$ with
$\rho_{i}=2^{-i(k-1)}\rho$, so $t_{i}\approx\lambda/(\log_{2}(1/\rho)+i(k-1))$. The sum over
$M=m/k$ rounds is harmonic, not linear; using $\sum_{i=1}^{M}1/(i+c)<\log(M/c+1)+1$ (STIR
Fact C.1),
\[
  q = O\!\left(\frac{\lambda}{\log(1/\rho)}
    + \frac{\lambda}{k}\log\!\Bigl(\frac{m}{k\log(1/\rho)}\Bigr)
    + \frac{m}{k}\right)
  \qquad\text{versus FRI's}\quad O\!\left(\lambda+\frac{\lambda m}{k}\right).
\]
In the typical regime $m\le\lambda$, $\rho=O(1)$, setting $k\approx\log m$ yields $q=O(\lambda)$,
which is optimal. \textbf{There is no further protocol-level lever on query count.}

\section{A STIR iteration: quotient and degree correction}

STIR sends a folded $g$ on a larger domain, pins it with an out-of-domain sample, checks it against
$t$ shift queries, and forms the next function as a quotient:
\[
  \mathrm{Quotient}(f,S,\mathrm{Ans},\mathrm{Fill})(x)
  =\frac{f(x)-\widehat{\mathrm{Ans}}(x)}{\hat V_{S}(x)}\qquad (x\notin S).
\]
Two of those symbols carry the whole mechanism. $S=\{z_{0},z_{1},\dots,z_{t}\}$ is the out-of-domain
sample together with the shift points; $\widehat{\mathrm{Ans}}$ interpolates the values the verifier
computed \emph{for itself} at those points; and
\[
  \hat V_{S}(X)=\prod_{z\in S}(X-z)
\]
is the \emph{vanishing polynomial} of $S$ --- monic, of degree $|S|=t+1$, zero exactly on $S$.

The quotient is a polynomial only when the numerator vanishes wherever the denominator does, so
\begin{align*}
  \frac{g-\widehat{\mathrm{Ans}}}{\hat V_{S}}\ \text{is a polynomial}
  \;&\Longleftrightarrow\;
  g-\widehat{\mathrm{Ans}}\ \text{vanishes on } S\\
  \;&\Longleftrightarrow\;
  g\ \text{carries the verifier's answers at } S.
\end{align*}
That equivalence \emph{is} the enforcement. Nobody checks the answers directly; the division either
comes out clean or it does not, and when it does not the result is not slightly wrong but far from
low degree (Lemma 4.4) --- so the next round's proximity test, which was running anyway, rejects it.
A lie is converted into a distance the existing machinery already measures.

Two consequences follow from $\hat V_{S}$ alone. The side condition $x\notin S$ is not fussiness:
the denominator vanishes there, which is why the new domain must be chosen disjoint from $S$ and why
$\mathrm{Fill}$ exists to supply values at the holes. And the degree drops by $|S|$, so the round
closes with degree correction --- a geometric sum evaluated in $O(\log e)$ operations by repeated
squaring rather than term by term.

\section{A WHIR iteration: constraints instead of quotients}

\begin{enumerate}
  \item \textbf{Sumcheck rounds.} $k$ rounds on
        $\sum_{\boldsymbol b}\hat w(\hat f(\boldsymbol b),\boldsymbol b)=\sigma$, giving
        $\hat h_{1..k}$ and randomness $\boldsymbol\alpha$.
  \item \textbf{Folded function.} The prover sends $g$ on $\mathcal{L}^{(2)}$.
  \item \textbf{Out-of-domain sample.} $z_{0}\leftarrow\mathbb{F}$, reply $y_{0}$.
  \item \textbf{Shift queries.} $z_{1},\dots,z_{t}\leftarrow\mathcal{L}^{(2^{k})}$, answers
        $y_{i}=\Fold(f,\boldsymbol\alpha)(z_{i})$; then a batching challenge $\gamma$.
  \item \textbf{Recursive claim.}
\end{enumerate}
\begin{align*}
  \hat w'(Z,\boldsymbol X) &:= \hat w(Z,\boldsymbol\alpha,\boldsymbol X)
    + Z\cdot\sum_{i=0}^{t}\gamma^{\,i+1}\cdot\eqp(\boldsymbol z_{i},\boldsymbol X),\\
  \sigma' &:= \hat h_{k}(\alpha_{k}) + \sum_{i=0}^{t}\gamma^{\,i+1}\cdot y_{i}.
\end{align*}
The out-of-domain answer and the $t$ shift answers become $t+1$ additional \emph{constraints},
folded into the weight polynomial by powers of $\gamma$. No quotient, hence no degree correction,
hence --- as the paper states explicitly --- \textbf{the verifier performs no divisions at all}. On
a machine without \texttt{OP\_DIV}, where an inversion must be hinted and checked, that is the
reason this project should target WHIR rather than STIR.

\section{What a shift query actually is}\label{sec:shift}

The most frequently occurring and most easily skimmed quantity in the protocol is
$r^{\mathrm{shift}}_{i}$ (WHIR writes $z_{i}$). It deserves its own section, because in this
implementation it \emph{is} ``a query'', and the identification is not obvious.

\subsection*{Which domain it lives in}

The first thing to get straight: a shift query is \textbf{not} drawn from the new domain
$\mathcal{L}'$ the prover just committed to. It is drawn from
\[
  r^{\mathrm{shift}}_{i}\leftarrow \mathcal{L}^{(2^{k})}=\{x^{2^{k}}:x\in\mathcal{L}\},
  \qquad |\mathcal{L}^{(2^{k})}|=|\mathcal{L}|/2^{k},
\]
the domain the \emph{folded} function lives on. And the folded function is \textbf{virtual} --- no
one sends it. To read it at a point the verifier must read the original $f$ across the whole fibre:
\[
  y_{i}=\Fold(f,\boldsymbol\alpha)(r^{\mathrm{shift}}_{i}),
  \qquad\text{requiring } f \text{ on } \{x: x^{2^{k}}=r^{\mathrm{shift}}_{i}\},
  \text{ which is } 2^{k} \text{ values}.
\]
Those $2^{k}$ values are the paper's ``alphabet $\mathbb{F}^{2^{k}}$'' --- one query reads one
symbol, and a symbol is $2^{k}$ field elements. In the implementation it is a Merkle \textbf{row},
and it is exactly the row that has to be bound to its leaf.

\subsection*{Why it has to exist}

This is the pivot of the whole STIR/WHIR construction, and the rate reduction of \S4 is what
it pays for. In FRI the folded function lives \emph{naturally} on $\mathcal{L}^{(k)}$ and the
verifier can compute it unaided. To lower the rate, STIR and WHIR have the prover re-commit a $g$ on
a \emph{larger} new domain.

But nothing forces $g$ to be $\Fold(f,\boldsymbol\alpha)$. The prover could send an unrelated
low-degree polynomial. \textbf{The shift queries are the pin that fastens the new oracle back to the
old one}: they require $\hat g(r^{\mathrm{shift}}_{i})=\Fold(f,\boldsymbol\alpha)(r^{\mathrm{shift}}_{i})$
at $t$ random points. Hence the name: STIR $=$ \emph{Shift To Improve Rate}. Out-of-domain sampling
does a different job --- it pins $g$ to a \emph{unique} codeword in the list; the shift queries then
check that this codeword is the fold. Neither suffices alone.

\subsection*{It is the soundness term}

$t$ is not arbitrary; it is precisely the exponent in the query term of \S\ref{sec:eps}. The argument is
two lines: if the codeword $u$ near $g$ disagrees with $\Fold(f,\boldsymbol\alpha)$ on at least a
$\delta$ fraction, each uniform $r^{\mathrm{shift}}_{i}$ independently lands on a disagreement with
probability $\ge\delta$, so
\[
  \Pr[\text{all } t \text{ points miss}]\le(1-\delta)^{t}
  \;=\;\text{the leading term of }\varepsilon^{\mathrm{shift}}.
\]
So $t_{i}=\lceil\lambda/-\log_{2}(1-\delta_{i})\rceil$, and rate reduction saves money precisely
because it raises the usable $\delta$ and so lowers $t$. \S4's table and this are two sides of
one fact.

\subsection*{What it is in this implementation}

\begin{center}\small
\begin{tabular}{@{}p{5.2cm}p{8.4cm}@{}}
\toprule
\textbf{In the paper} & \textbf{In the script}\\
\midrule
sampling $r^{\mathrm{shift}}_{i}$ & \texttt{challenger::sample(slot)}, one base element from the rate\\
its index in $\mathcal{L}^{(2^{k})}$ & \texttt{field::low\_bits\_to\_altstack} --- \textbf{these bits are the Merkle path}\\
the $2^{k}$ values on the fibre & the opened row, bound to its leaf by \texttt{merkle::hash\_row}\\
$y_{i}=\Fold(\cdot)(r^{\mathrm{shift}}_{i})$ & \texttt{multilinear::eval\_multilinear(k)} folds the row\\
batching into $\sigma'$ & \texttt{constraint::combine\_answers(t)}\\
the accumulated pairs & \texttt{constraint::closing\_check}\\
\bottomrule
\end{tabular}
\end{center}

The second row is worth pausing on. \textbf{The index is the path} --- the sampled bits are directly
the direction bits of the Merkle walk, so the spender picks neither which leaf is opened nor where
it sits; the binding is structural rather than an added check. It also explains where the
$\mathrm{depth}$ factor in the cost law comes from:
$\mathrm{depth}=\log_{2}|\mathcal{L}^{(2^{k})}|$, so one shift query is $\mathrm{depth}$ Poseidon2
permutations, and that is the origin of the 1300 in \S\ref{sec:cost}.

\section{Soundness, and what it rests on}\label{sec:eps}

Compilation through BCS requires \emph{round-by-round} soundness. WHIR's Theorem 5.2 gives:
\begin{align*}
  \varepsilon^{\mathrm{fold}}_{i,s} &\le \frac{d\cdot\ell_{i,s-1}}{|\mathbb{F}|}
    + \err^{\star}(\mathcal{C},2,\delta_{i}),
  &\varepsilon^{\mathrm{out}}_{i} &\le \frac{2^{m_{i}}\ell^{2}_{i,0}}{2|\mathbb{F}|},\\
  \varepsilon^{\mathrm{shift}}_{i} &\le (1-\delta_{i-1})^{t_{i-1}}
    + \frac{\ell_{i,0}(t_{i-1}+1)}{|\mathbb{F}|},
  &\varepsilon^{\mathrm{fin}} &\le (1-\delta_{M-1})^{t_{M-1}}.
\end{align*}
Only $\varepsilon^{\mathrm{shift}}$ and $\varepsilon^{\mathrm{fin}}$ carry the $(1-\delta)^{t}$
term; everything else is field-size noise.

Because WHIR discards quotienting, distance preservation under folding is not enough: the analysis
needs \emph{list} preservation, which is the paper's new tool.

\medskip
\noindent\fbox{\begin{minipage}{0.96\linewidth}\small
\textbf{Definition 4.9 (mutual correlated agreement, informal).} $\mathsf{Gen}$ has mutual
correlated agreement if, with high probability over $\alpha$: for \emph{every} set $S$ on which
$\sum_{j}r_{j}f_{j}$ agrees with $\mathcal{C}$, \emph{every} $f_{i}$ agrees with $\mathcal{C}$ on
$S$.
\end{minipage}}
\medskip

Standard correlated agreement promises only \emph{some} set $T$; equating $S$ and $T$ upgrades
``folding preserves distance'' to ``every codeword close to $\Fold(f,\alpha)$ is the folding of a
codeword close to $f$'' (Theorem 4.20). It is proved in the unique-decoding regime and
\emph{conjectured} beyond it (Conjecture 4.12).

\begin{center}\small
\begin{tabular}{@{}llcrr@{}}
\toprule
\textbf{Regime} & \textbf{Proximity }$\delta$ & \textbf{Status} & \textbf{Verifier hashes} & \textbf{Argument}\\
\midrule
WHIR-UD & $(1-\rho)/2$            & \ok                   & 10\,000 & 621 KiB\\
WHIR-JB & $1-\sqrt{\rho}-\eta$    & \warn{Conj.\ 4.12(1)} & 5\,100 & 299 KiB\\
WHIR-CB & $1-\rho-\eta$           & \bad{Conj.\ 4.12(2)}  & 2\,700 & 156 KiB\\
\bottomrule
\end{tabular}
\end{center}
\begin{center}\footnotesize\color{muted}
WHIR Table 5, $m=24$, $\rho=1/2$, $\lambda=128$, Goldilocks cubic extension.
\end{center}

\clearpage
\begin{center}{\large\bfseries Part II\quad Implementation review and repair}\end{center}
\vspace{-0.5em}
\noindent{\color{rule}\rule{\linewidth}{0.6pt}}

\section{What the crate is, and its cost law}\label{sec:cost}

Two crates: \texttt{poseidon2} (the permutation, the quartic extension, Merkle verification over
KoalaBear) and \texttt{whir} (sponge, challenger, sumcheck, multilinear evaluation, constraint,
composed verifier). The design premise is sound and well argued --- an algebraic hash makes a digest
eight field elements, so compressing two children is sixteen stack items handed to an arithmetic
routine, and nothing needs concatenating. That removes the \texttt{OP\_CAT} soft-fork dependency
entirely.

Everything else follows from one measured constant: a Poseidon2 permutation is
\textbf{572\,228 bytes} of script, and nothing else is within 1\% of it. So
\[
  C_{\mathrm{script}}\approx N_{\mathrm{perm}}\cdot c_{\mathrm{perm}},
  \qquad
  N_{\mathrm{perm}}=\underbrace{\sum_{i}t_{i}\cdot\mathrm{depth}_{i}}_{\text{Merkle}}
  \;+\; N_{\mathrm{transcript}} .
\]
For the crate's example configuration the braced term is
$20\cdot(21{+}17{+}13{+}9)+20\cdot5=\mathbf{1300}$ permutations. \textbf{Every optimisation that
matters is an optimisation of that term.}

\section{Verified, versus asserted}

A test comparing a script against a Rust transcription establishes that the two \emph{agree}; only a
test driven by a real proof, with a negative case, establishes that either is \emph{right}.

\begin{center}\small
\begin{tabular}{@{}p{4.4cm}p{3.6cm}p{3.2cm}cc@{}}
\toprule
\textbf{Component} & \textbf{Evidence} & \textbf{Negative case} & \textbf{Before} & \textbf{Now}\\
\midrule
Poseidon2, Merkle compression & vs.\ reference & --- & \ok & \ok\\
Duplex sponge, padding & vs.\ reference & short absorb $\neq$ padded & \ok & \ok\\
Challenger sampling & vs.\ reference & coefficient order & \ok & \ok\\
Merkle opening of a real proof & real commitment root & tampered leaf rejected & \ok & \ok\\
Challenge derivation & squeezed in script & perturbed evaluation rejected & \bad{D1} & \fixed\\
Closing identity & EF identity & each coefficient rejected & \bad{D2} & \fixed\\
Composed \texttt{verify()} & executed vs.\ reference & transcript moves the claim & \bad{D3} & \fixed\\
Row bound to leaf & \texttt{hash\_row} vs.\ Plonky3 & every substitution rejected & \bad{gap} & \fixed\\
Answers batched into target & vs.\ reference & every answer moves $\sigma'$ & \bad{gap} & \fixed\\
Accumulated pairs threaded & --- & --- & \bad{gap} & \bad{open}\\
\bottomrule
\end{tabular}
\end{center}

\section{The three defects}

\defect{D1}{Fiat--Shamir challenges were never derived from the transcript}
The challenges fed to the sumcheck chain and the point at which the final polynomial was evaluated
were deterministic constants, invented by the test rather than derived by Plonky3.

Consider one round, with $h(1)$ recovered from $h(0)+h(1)=c$:
\[
  h(r)=(1-r)\,c_{0}+r\,(c-c_{0})+r(r-1)\,c_{\infty}.
\]
Fix any target $c'$ and any $r\notin\{0,1\}$. Since $r(r-1)\neq0$ is invertible,
\[
  c_{\infty}:=\bigl(c'-(1-r)c_{0}-r(c-c_{0})\bigr)\cdot\bigl(r(r-1)\bigr)^{-1}
\]
reaches $c'$ for \emph{any} $c_{0}$. Composing, a prover supplying the challenges controls the
closing claim exactly. \textbf{The soundness error of the chain was 1.}

\textbf{Fix.} \texttt{sumcheck\_round\_fs} absorbs the round's $(c_{0},c_{\infty})$ and squeezes $r$
from the resulting rate, so $r$ is a function of the proof and is fixed only after both evaluations
are committed. The invariant is \texttt{claim(4) || state(16)} with the sponge on top, already the
layout \texttt{absorb} wants, so rounds compose.
\defectend

\defect{D2}{The closing identity was a tautology, in the base field}
The test pushed $a$, $b$ and $ab$, rotated twice, and \texttt{final\_check} recomputed $ab$ to
compare against $ab$ --- true by construction for any proof. It also compared one coefficient,
leaving a base-field equality worth $\le\log_{2}P\approx31$ bits where the quartic extension exists
to provide $\approx$124.

\textbf{Fix.} \texttt{final\_check} multiplies in $\mathrm{EF}$ and compares all four coefficients;
the end-to-end weight is copied off the stack rather than pushed, tying the sumcheck chain to the
final polynomial.
\defectend

\defect{D3}{\texttt{verify()} priced a schedule it could not run}
\texttt{sample\_ef} was never called in \texttt{src/}; \texttt{absorb} leaves a sixteen-element state
where \texttt{sumcheck\_round} expects \texttt{claim c0 c\_inf r}; and the test only took
\texttt{.len()}. Two bugs went with it: \texttt{open\_query} hardcoded rate slot 0, so every query in
a round derived the \emph{same} index and opened the \emph{same} leaf; and a round absorbed two
elements where it sends two \emph{extension} elements, then paid a second permutation for a squeeze
that absorbing had already performed.

\textbf{Fix.} Wired, executed against a reference, with the slot parameterised and the accounting
corrected.
\defectend

\section{The repair's stack discipline}

\begin{lstlisting}
// in : claim(4) || state(16)        hints pushed above: c0(4) || c_inf(4)
// out: claim'(4) || state'(16)
{ ext4::copy(0) } { ext4::to_altstack() }   // park c_inf -- absorb consumes it
{ ext4::copy(1) } { ext4::to_altstack() }   // park c0
{ sponge::absorb(8) }                       // ONE permutation; absorb IS the duplexing
{ challenger::sample_ef(0) }                // r := rate[0..4], bound to the transcript
for _ in 0..4  { { 23 } OP_ROLL }           // lift claim over state'  (2*D + WIDTH - 1)
{ ext4::from_altstack() }                   // c0
{ ext4::from_altstack() }                   // c_inf
for _ in 0..4  { { 15 } OP_ROLL }           // lift r to the top       (4*D - 1)
{ sumcheck::sumcheck_round() }              // consumes claim c0 c_inf r
for _ in 0..16 { { 19 } OP_ROLL }           // restore the invariant   (D + WIDTH - 1)
\end{lstlisting}

Each \texttt{OP\_ROLL} group lifts the deepest slot of its element first, preserving coefficient
order.

\medskip
\noindent\fbox{\begin{minipage}{0.96\linewidth}\small
These depths are the failure mode of the whole approach.
\texttt{CLAIM\_UNDER\_STATE\_AND\_R} is $2D+W-1=23$ and not $4D+W-1=31$; the difference between
those two is one misread of what sits above the claim, and it produces
\texttt{InvalidStackOperation} rather than a wrong answer. That is a class of error which appears
only when the script is \emph{executed}, which is a footnote to D3: a builder that is measured and
never run cannot be wrong in this way, because it is never right in it either.
\end{minipage}}

\section{What the repair cost}

Rolls and picks are one to two bytes against 572\,228 for a permutation, so the wiring is free at
this resolution, and the round drops from two permutations to one.

\begin{center}\small
\begin{tabular}{@{}lrrrp{4.8cm}@{}}
\toprule
\textbf{Quantity} & \textbf{Before} & \textbf{Predicted} & \textbf{Measured} & \textbf{Why}\\
\midrule
Transcript permutations & 67 & $\approx$46 & \textbf{44} & absorb is the duplexing\\
Merkle paths            & 1300 & 1300 & 1300 & the term that matters\\
Transcript share        & 5.0\% & $\approx$3.5\% & \textbf{3.3\%} & ---\\
Tests                   & --- & --- & \textbf{67} & all passing\\
\bottomrule
\end{tabular}
\end{center}

The honest verifier is smaller than the priced skeleton. \textbf{There is no efficiency argument for
leaving Fiat--Shamir out.} The composed end-to-end script measures \textbf{8\,967\,342 bytes}: a
real Plonky3 WHIR proof accepted by one Bitcoin Script execution, with every perturbation of the
prover's evaluations rejected.

\section{Constraining the query}\label{sec:constrain}

Three changes, in the order they had to happen.

\subsection*{The row is bound to its leaf}
A query opening has two halves and only one used to be checked. Walking the path proves that
\emph{some} committed leaf sits at the sampled index; with the leaf arriving as a hint alongside the
row, nothing stopped a spender authenticating the real leaf and folding a different row.
\texttt{merkle::hash\_row} reproduces Plonky3's \texttt{PaddingFreeSponge}: overwrite the first
$\mathrm{RATE}$ state elements with each chunk and permute, a partial final chunk leaving the rest
alone. No padding, no length tag --- deliberately \emph{not} the duplex, where zeroing the rest of
the rate and binding the count is what stops a short absorb colliding with a padded longer one.
Cost: $\lceil \mathrm{row\_len}/8\rceil$ permutations per query.

\subsection*{A round's answers are batched}
\texttt{constraint::combine\_answers} computes $\sigma'=\hat h_{k}(\alpha_{k})+\sum_{i}\gamma^{i+1}y_{i}$
by Horner from the top --- $t+1$ multiplications rather than the $2t$ of forming the powers, the
outer $\gamma$ being what shifts every exponent up by one. No permutation at all.

\subsection*{The closing identity, and two readings of it}
Each accumulated term collapses, because summing an evaluation against $\eqp$ \emph{is} evaluation:
\[
  \sum_{\boldsymbol b} \hat f_{M}(\boldsymbol b)\,\gamma^{e}\,\eqp(\boldsymbol z,\boldsymbol b)
  = \gamma^{e}\,\hat f_{M}(\boldsymbol z).
\]
So $\hat w'$ is exactly the list of $(c_{j},\boldsymbol z_{j})$ pairs of \S\ref{sec:crs}, and checking it means
evaluating the final polynomial at each point and taking the weighted sum. That is self-consistent,
it is the paper's reading, and \texttt{closing\_check} implements it.

It is not the identity Plonky3 checks. There the two polynomials are kept \emph{apart}, each
evaluated at its own point:
\[
  \underbrace{\texttt{claimed\_eval}}_{\text{the chained claim}}
  \;=\;
  \underbrace{w(\boldsymbol R)}_{\texttt{eval\_constraints\_poly}}
  \;\cdot\;
  \underbrace{\hat f_{M}(\boldsymbol r_{\mathrm{fin}})}_{\texttt{final\_evaluations.eval\_ext}},
\]
and \texttt{eval\_constraints\_poly} evaluates every accumulated $\eqp(\boldsymbol z_{j},\cdot)$
explicitly. Since this implementation is checked against Plonky3, that is the shape the script must
produce --- so an $\eqp$ primitive is needed after all.

The two readings are easy to mistake for one, and the mistake runs in a particular direction: the
collapse above is a true statement about the paper, and reading it as a statement about the verifier
being targeted removes a primitive that verifier requires. The gap appears only when the two
formulations are set side by side.

Measured, since this is where \texttt{final\_poly\_vars} bites:

\begin{center}\small
\begin{tabular}{@{}rr@{}}
\toprule
$n_{\mathrm{vars}}$ & bytes per point\\
\midrule
2 & 95{,}392\\
3 & 190{,}922\\
4 & 381{,}988\\
5 & 764{,}270\\
\bottomrule
\end{tabular}
\end{center}

Doubling per variable. For the example configuration the accumulated pairs number roughly
$\sum_i (t_i+2)\approx110$, so at $\texttt{final\_poly\_vars}=4$ the closing check is on the order of
42\,MB against 743\,MB of Merkle work --- about 5\%. It is also a concrete argument for pushing
folding work into earlier rounds: the cost is $n_{\mathrm{points}}\cdot2^{n_{\mathrm{vars}}}$.

\subsection*{Threading, and why it costs almost nothing}

The list of pairs has to reach the closing check having been \emph{derived}, not supplied. The
natural guess is that this is expensive: a pair created early should have to be rewritten every
round, its coefficient absorbing a factor as the sumcheck binds variables. It does not, and the
reason is one line --- $\eqp$ factorises across a split,
\[
  \eqp\bigl((\boldsymbol z_{\mathrm{pre}},\boldsymbol z_{\mathrm{suf}}),
            (\boldsymbol\alpha,\boldsymbol y)\bigr)
  =\eqp(\boldsymbol z_{\mathrm{pre}},\boldsymbol\alpha)\cdot
   \eqp(\boldsymbol z_{\mathrm{suf}},\boldsymbol y),
\]
so binding a prefix never mixes coordinates. Constraints can simply be \emph{accumulated}, and at
the end each is evaluated against the concatenation of every round's folding randomness,
\[
  \boldsymbol R=\boldsymbol\alpha^{(0)}\,\|\,\boldsymbol\alpha^{(1)}\,\|\cdots\|\,
  \boldsymbol\alpha^{(\mathrm{fin})},
\]
a constraint of arity $n_{c}$ reading the \textbf{last $n_{c}$ coordinates} of $\boldsymbol R$ ---
which is what it means for the randomness from round $i$ onwards to be what binds a constraint born
at round $i$.

On a stack that rule is free. Push $\boldsymbol R$ with $R_{0}$ deepest and the last $n_{c}$
coordinates are the \emph{topmost} $n_{c}$: every constraint reads from the top, and a longer one
reaches further down. No slicing, no offset arithmetic, no per-round rewriting. Keeping each round's
challenge rather than discarding it costs \textbf{76 bytes}, 0.011\% of a round.

A second convenience sits in the point expansion. An out-of-domain sample is a scalar $u$ whose
constraint point is $(u^{2^{m-1}},\dots,u^{2},u)$; repeated squaring produces $u$ first, so moving
the block to the altstack top-first lands the coordinates in exactly the order the $\eqp$ evaluation
consumes them. The two conventions meet without a roll.

\begin{center}\small
\begin{tabular}{@{}llr@{}}
\toprule
\textbf{Primitive} & \textbf{What it is} & \textbf{Cost}\\
\midrule
\texttt{eq\_eval} & $\prod_i (1+2z_ir_i-z_i-r_i)$ & $2n$ ext.\ mult.\\
\texttt{expand\_univariate} & the square-power point & $n-1$ ext.\ mult.\\
\texttt{constraint\_eval} & $w(\boldsymbol R)$ & per constraint, above\\
\texttt{sumcheck\_rounds\_fs\_keep} & buries each round's $\alpha$ & 76 B per round\\
\texttt{constraint\_eval\_batched} & the same, from scalars and $\chi$ & $+32$--$45\%$\\
\bottomrule
\end{tabular}
\end{center}
\begin{center}\footnotesize\color{muted}
None costs a permutation. \texttt{eq\_eval} and \texttt{expand\_univariate} are checked against
Plonky3's own \texttt{Point::eval\_eq} and \texttt{Point::expand\_from\_univariate}, not only against
a local reference --- either could otherwise have been a consistent misreading.
\end{center}

\subsection*{Deriving the constraints, and a convention worth pinning}

The list can be \emph{derived} rather than supplied, and far more cheaply than carrying it whole: a
round's constraints are all weighted by powers of one batching challenge, and each of their points is
the square-power expansion of a single scalar, so a round is $1+n$ extension elements rather than
$n(1+\text{arity})$.

The weights needed pinning down. Plonky3's \texttt{Constraint::challenge\_powers(shift)} is
\texttt{shifted\_powers}$(\chi^{\text{shift}})$, so within a constraint the weights run
$\chi^{0},\chi^{1},\chi^{2},\dots$ --- starting at \textbf{one}. The $\sigma'$ update in the same
protocol uses $\gamma^{\,i+1}$, starting at the challenge itself. \textbf{Two conventions inside one
protocol} is exactly the sort of thing a script and a reference written by the same hand get wrong
together, so the two-constraint case is pinned against the expression written out by hand rather than
against either implementation. Starting at one also makes a group a plain polynomial in $\chi$, so
Horner evaluates it in $n-1$ multiplications rather than the $2n$ a power ladder costs.

\subsection*{An absorb four times too narrow}

Wiring the threading in exposed a defect the earlier accounting had hidden. Out-of-domain answers
and the final polynomial are \emph{extension} elements, and Plonky3 observes them with
\texttt{observe\_algebra\_slice} --- four base elements each, not one. The out-of-domain points are
also sampled one at a time \emph{between} the answers rather than once after them, so a round costs
one duplexing per sample and not one per rate block. Both are the kind of error that an accounting
cannot catch and an execution can: nothing about the schedule looks wrong until something absorbs a
value and the challenge that follows it comes out different.

\section{Composing the query openings}\label{sec:compose}

The queries were priced and never emitted, and emitting them found two defects and one omission.

\defect{D4}{\texttt{open\_query} could not have run}
It took the opening from the main stack and sampled the index with a pick that reaches into the top
sixteen slots --- so with an opening above the sponge, those picks landed inside the row. The
contract could not hold, and nothing caught it because the builder \textbf{had no test at all}. It is
the same shape as D3: a script that is measured and never executed.

\textbf{Fix.} The opening comes off the altstack and the index is sampled with a pick that reaches
past it. The order is forced rather than chosen: the direction bits go on the altstack, so the
opening has to come off first or the bits would be popped in its place.
\defectend

\defect{D5}{The opening's root came from the witness}
A spender could commit to one tree and open a path in another. Any root with a consistent path
passed, and the round's \emph{absorbed} commitment --- the one the transcript is bound to --- was
never compared against.

\textbf{Fix.} The root is copied from that commitment, which the absorb now keeps. This is the third
value in the crate found to be absorbed and then not retained, after the final polynomial and the
sumcheck randomness; the pattern is worth naming, because in each case the transcript proved the
prover had \emph{said} something and nothing tied what was said to what was later used.
\defectend

Keeping it needed care about \emph{where}. A sumcheck round buries its challenge directly under the
claim, so a commitment parked there would land inside $\boldsymbol R$ and stop the suffix rule of
\S\ref{sec:constrain} working. Commitments go below $\boldsymbol R$, where they survive the round's sumcheck rounds
--- which they must, because Plonky3 opens round $i$'s queries against round $i-1$'s commitment, not
the one round $i$ just absorbed. Two are alive during a round, the older deeper by exactly one digest.

\subsection*{The omission: the leaf hash was never priced}

The row-to-leaf binding of \S\ref{sec:constrain} closed a real gap, but the accounting only ever counted the path. At
folding factor 4 a row is sixteen extension elements, so sixty-four base elements, so \textbf{eight
permutations per query} on top of the twenty-one the path costs.

\begin{center}\small
\begin{tabular}{@{}lrr@{}}
\toprule
\textbf{Example configuration} & \textbf{before} & \textbf{now}\\
\midrule
Merkle paths & 1300 & 1300\\
leaf hashing & \emph{unpriced} & \textbf{800}\\
transcript & 54 & 62\\
\midrule
total & 1354 & \textbf{2162}\\
query share & --- & \textbf{97.1\%}\\
\bottomrule
\end{tabular}
\end{center}

\subsection*{Deriving the constraints, and what it costs}

The closing check was sound for the statement and merely \emph{described} everything else: its
constraint list came from the witness, so the out-of-domain and shift answers were authenticated and
then ignored. They are now derived, and what makes that affordable is that a constraint is far
smaller than it looks. A round's constraints share one batching challenge and each of their points is
the square-power expansion of a single scalar, so a round carries $1+n$ extension elements rather
than $n(1+\text{arity})$. The rounds bury those scalars below the randomness as they sample them and
the closing check lifts them back out.

Three things had to line up. \textbf{The order}: Horner wants a group's last term first, and a round
produces its scalars in order and its challenge afterwards --- so the $\eqp$ values are built on the
stack as the scalars arrive and Horner runs once the challenge appears. \textbf{The shift scalars}:
a shift query's point is $g^{\,\text{index}}$, and the index is re-sampled rather than kept, because
sampling is a pick and the sponge has not moved, while keeping it would mean carrying it past a
Merkle walk that uses the altstack. \textbf{The arities}: a constraint born at round $i$ is bound by
the randomness from round $i$ onwards, so it reads the last $m-r_{\text{len}}$ coordinates ---
\S\ref{sec:constrain}'s suffix rule, now with the arities computed rather than passed in.

\begin{center}\small
\begin{tabular}{@{}lrr@{}}
\toprule
\textbf{Closing check} & \textbf{supplied} & \textbf{derived}\\
\midrule
bytes & 1{,}550{,}899 & 92{,}801{,}911\\
share of the verifier & 0.126\% & \textbf{7.5\%}\\
constraints evaluated & 1 & \textbf{110}\\
permutations & 0 & \textbf{0}\\
\bottomrule
\end{tabular}
\end{center}

Sixty times larger, and still not a permutation. Everything this report has said about the algebra
being cheap survives the composition; what grew, every time, is the hashing.

\medskip
\noindent\fbox{\begin{minipage}{0.96\linewidth}\small
\textbf{What is still supplied, and why it always will be.} The statement's own constraint. Its
point is public because it \emph{is} the claim being proved, and a different point is a different
statement rather than a cheaper proof of the same one. Two fidelity items remain against Plonky3:
each group ends with a fresh squeeze for its batching challenge rather than reading whichever rate
slot the query indices happened to leave, and Plonky3 folds the statement and the commit phase's
out-of-domain samples into one \texttt{initial\_constraint} where this treats them as two groups.
\end{minipage}}

\section{The stack the verifier keeps}\label{sec:stack}

Everything above is about what the verifier computes. This is about where it puts things, which on a
stack machine is most of the work and all of the risk: every defect in \S\ref{sec:compose} was a
value in the wrong place rather than an arithmetic mistake.

\subsection*{The invariant}

Between steps the main stack reads, bottom to top:
\[
  \underbrace{\text{carried scalars}}_{\text{constraints}}
  \;\|\;
  \underbrace{\hat f_{M}}_{2^{v}\text{ elements}}
  \;\|\;
  \underbrace{\boldsymbol R}_{m\text{ elements}}
  \;\|\;
  \underbrace{\text{claim}}_{1}
  \;\|\;
  \underbrace{\text{sponge}}_{16\text{ slots}} .
\]
Three constraints fix that order, and between them there is no freedom left.

\smallskip
\noindent\textbf{The sponge is on top} because absorbing wants its inputs directly above the state
and sampling is a pick into the top sixteen slots. Everything else has to be underneath, which is why
a query opening --- the one thing that must sit above the state while it is checked --- needed a
sampling primitive that reaches past it.

\smallskip
\noindent\textbf{$\boldsymbol R$ must be contiguous}, because \S\ref{sec:constrain}'s rule reads a
constraint's local point as a \emph{suffix} of it. A sumcheck round appends its challenge directly
under the claim, so anything else parked there would land in the middle of the randomness and the
suffix rule would silently read the wrong coordinates. That single requirement is what forces
commitments and the final polynomial below $\boldsymbol R$ rather than under the claim, and it is not
a preference: nothing would fail loudly if it were violated.

\smallskip
\noindent\textbf{The altstack is the transcript.} Everything the spender supplies is on it in
consumption order, so nothing may be parked there across a step. A value needing to survive one has
to be buried on the main stack instead --- which is the whole reason the layout above exists.

\subsection*{Why every depth is a constant}

Burying a value under a region, and copying one back out from under it, are both loops at a
\emph{fixed} depth. The argument is the same in both directions and worth stating once, because it is
used a dozen times and looks like an off-by-one every time:

\begin{quote}
Rolling the deepest slot of a region to the top removes one item from below the cursor and adds one
above it. The two cancel, so the next-deepest slot is at the same depth as the last one was.
\end{quote}

So ``lift a region of $k$ slots over a block of $b$'' is $k$ rolls at depth $b+k-1$, and it preserves
the region's internal order. The same cancellation is why \texttt{ext4::copy} is four picks at one
depth, and why copying a block back out is $b$ picks at the depth the region left it.

The one place this does \emph{not} hold is dropping from the bottom: each slot removed shrinks the
stack, so the depth decreases by one per step. Both forms appear, and confusing them is exactly the
class of error that only shows up when the script is executed.

\subsection*{What a round does to the layout}

\begin{center}\small
\begin{tabular}{@{}lll@{}}
\toprule
\textbf{Step} & \textbf{Effect on the layout} & \textbf{Depth} \\
\midrule
absorb the commitment & one segment below $\boldsymbol R$ & $4r+20+8-1$\\
out-of-domain sample & one element below $\boldsymbol R$ & $4r+20+4-1$\\
query opening & nothing; balanced & reaches the commitment\\
shift scalar & one element below $\boldsymbol R$ & $4r+20+4-1$\\
batching challenge & one element below $\boldsymbol R$ & $4r+20+4-1$\\
drop the opened commitment & a segment leaves the middle & recomputed\\
sumcheck rounds & $\boldsymbol R$ grows under the claim & $2\cdot4+16-1$\\
\bottomrule
\end{tabular}
\end{center}

Only the second-to-last row is awkward. Dropping a commitment removes a segment from the
\emph{middle} of the carried region, so every depth above it shifts --- which is why the generator
keeps a model of the region rather than hard-coding the numbers. Two commitments are live at once,
because a round's queries open the \emph{previous} round's codeword while the current one has already
been absorbed.

\subsection*{And what the closing check does with it}

In order: drop the sponge; lift $\hat f_{M}$ over $\boldsymbol R$ and the claim and collapse it to one
value; evaluate the statement's constraint, which \emph{reads} $\boldsymbol R$ without consuming it;
lift the carried scalars out and hand them to the altstack; evaluate the accumulated constraints,
which consumes $\boldsymbol R$; and compare.

The order is not free. The statement has to be evaluated before the scalars reach the altstack,
because its pairs are down there underneath them; and it has to be evaluated before the accumulated
constraints, because those consume the randomness it needs. Moving the scalars reverses their order
twice --- once rolling them up, once moving them across --- which is what lets a round bury them in
the order it samples them and still have Horner read them in the order it wants.

\section{Two unused levers}\label{sec:levers}

\subsection*{Proof of work}
\texttt{pow\_bits: 0} is why the security level tops out at 100. Both reference implementations
grind: STIR's benchmarks take 22 bits, reducing the soundness required from the IOP from $2^{-128}$
to $2^{-106}$; WHIR uses $m-\log\rho-3$, also 22 at $m=24$, $\rho=1/2$. The trade is unusually
favourable here --- a query costs $\mathrm{depth}\times572$\,KB, a proof-of-work check costs one
permutation regardless of depth:
\[
  \Delta N_{\mathrm{perm}} = -\lceil20/4\rceil\cdot9 + 1 = -44 \ \text{of}\ 234
  \;\approx\; \textbf{19\% of the dominant term}.
\]

\subsection*{Soundness regime}
\texttt{SecurityAssumption::CapacityBound} is the cheapest of the three and the only one resting on
the unproven half of Conjecture 4.12. Nothing in the implementation mentions a soundness assumption:
the script executes a schedule, and the regime reaches it only as a query count. So the choice is a
configuration change, and \texttt{whir/tests/budget.rs} prices it --- each regime given the grinding
it needs, since holding \texttt{pow\_bits} fixed reports \emph{not derivable} for the weaker
assumptions rather than a price.

\begin{center}\small
\begin{tabular}{@{}lrrrrr@{}}
\toprule
\textbf{Regime} & $m$ & \texttt{pow} & \textbf{queries} & \textbf{permutations} & \textbf{vs CB}\\
\midrule
UD, proved            & 16 & 0  & 110 & 7\,376  & \textbf{7.18$\times$}\\
JB, Conj.\ 4.12(1)    & 16 & 26 & 39  & 1\,853  & 1.80$\times$\\
CB, Conj.\ 4.12(2)    & 16 & 19 & 21  & 1\,027  & 1.00$\times$\\
\midrule
UD, proved            & 20 & 0  & 110 & 11\,171 & \textbf{9.01$\times$}\\
JB, Conj.\ 4.12(1)    & 20 & 33 & 35  & 2\,222  & 1.79$\times$\\
CB, Conj.\ 4.12(2)    & 20 & 28 & 19  & 1\,240  & 1.00$\times$\\
\bottomrule
\end{tabular}
\end{center}

Two things the paper's own comparison does not show at this instance. \textbf{UD wants no grinding at
all}, paying entirely in queries, and its count is capped at 110 regardless of size --- so the gap
widens with $m$ rather than staying fixed. And \textbf{JB wants more grinding than CB}, 26 bits
against 19, which is the direction a weaker assumption should push: less can be assumed, so more has
to be bought.

A defensible choice, and the papers' own benchmarking default. But the proved alternative costs 7 to
9 times here, not the 3.7 of WHIR's Table 5 --- that table is Goldilocks with a cubic extension at
$m=24$ and rate $1/2$, and a ratio does not travel between instances. A verifier that settles bitcoin
should say out loud which assumption it stands on, and what the proved one would cost \emph{here}.

\clearpage
\begin{center}{\large\bfseries Part III\quad Completeness and soundness}\end{center}
\vspace{-0.5em}
\noindent{\color{rule}\rule{\linewidth}{0.6pt}}

\medskip
This part separates two things: the completeness and soundness of the \textbf{protocol} (the papers'
results, stated here and explained in shape), and of \textbf{this script} (this repository's
responsibility, covering only what it actually checks). Conflating them is the error this report has
repeatedly pointed at.

\section{Completeness}

\textbf{Proposition (protocol).} If $f\in\CRS[\mathbb{F},\mathcal{L},m,\hat w,\sigma]$ and the prover
is honest, the verifier accepts with probability \textbf{1}.

\emph{Proof sketch.} Each equality is an identity in the honest case, not a probabilistic event.

\begin{enumerate}
  \item \textbf{Sumcheck.} The honest $\hat h_{\ell}(X)=\sum_{\boldsymbol b}
  \hat w(\hat f(\boldsymbol\alpha,X,\boldsymbol b),\dots)$ makes
  $\hat h_{1}(0)+\hat h_{1}(1)=\sigma$ and
  $\hat h_{\ell}(0)+\hat h_{\ell}(1)=\hat h_{\ell-1}(\alpha_{\ell-1})$ hold by the definition of the
  sum.
  \item \textbf{Folding preserves codewords.} By Claim 4.15, $f\in\RS$ implies
  $\Fold(f,\boldsymbol\alpha)\in\RS$ with multilinear extension $\hat f(\boldsymbol\alpha,\cdot)$, so
  the honest $g$ is a codeword of the right degree on the right domain.
  \item \textbf{Answers.} The honest prover sets $y_{0}=\hat g(\boldsymbol z_{0})$, and $y_{i}$ is
  the $\Fold$ value the verifier computes from $f$ itself; the two agree by construction.
  \item \textbf{Recursive claim.} Substituting (2) and (3) into the definitions of $\hat w'$ and
  $\sigma'$ gives $\sum_{\boldsymbol b}\hat w'(\hat g(\boldsymbol b),\boldsymbol b)=\sigma'$, so the
  inductive hypothesis is preserved.
  \item \textbf{Closing.} $\hat f_{M}$ is sent in the clear and both final checks hold by
  construction. \hfill$\square$
\end{enumerate}

\textbf{Proposition (script).} Given a proof Plonky3 accepts and a witness laid out as documented,
the script runs to completion without aborting.

\emph{Why this is measured, not argued.} \texttt{whir/tests/end\_to\_end.rs} produces a real proof
with Plonky3's prover, confirms it with Plonky3's \emph{own} verifier, and only then hands the same
object to the script. The order is deliberate: comparing a script against a Rust reference alone
would let both pass while misreading a proof in the same way.

\section{Soundness}

\subsection*{Protocol: round by round}

BCS/Fiat--Shamir needs \emph{round-by-round} soundness, not plain soundness: a state function
$\mathsf{State}(\mathrm{tr})\in\{0,1\}$ that is 1 on the empty transcript iff the statement is true,
that a prover cannot flip back once 0, that a verifier's next challenge flips with probability
$\le\varepsilon_{i}$, and whose being 0 on a full transcript implies rejection. The strengthening is
necessary because Fiat--Shamir lets a prover retry challenges, under which plain soundness does not
survive.

The four errors of \S\ref{sec:eps} split in two, and the split is useful to an implementer:

\medskip
\noindent\fbox{\begin{minipage}{0.96\linewidth}\small
\textbf{Field-size terms} ($\varepsilon^{\mathrm{fold}}$, $\varepsilon^{\mathrm{out}}$, and the
second term of $\varepsilon^{\mathrm{shift}}$): of the form
$\mathrm{poly}(2^{m},1/\rho)/|\mathbb{F}|$, shrunk by \emph{extending the field}.\\
\textbf{Query terms} (the leading $\varepsilon^{\mathrm{shift}}$ and $\varepsilon^{\mathrm{fin}}$):
$(1-\delta)^{t}$, shrunk only by \emph{querying more}, and on Bitcoin each query is $\mathrm{depth}$
permutations.
\end{minipage}}
\medskip

The whole cost structure sits on that line: field size is free (the quartic extension of KoalaBear
gives about 124 bits), query count is the only thing that is paid for, and rate reduction exists to
lower it.

\subsection*{Script: what it actually guarantees}

\begin{center}\small
\begin{tabular}{@{}p{3.9cm}p{7.2cm}r@{}}
\toprule
\textbf{Property checked} & \textbf{Ground} & \textbf{Error}\\
\midrule
Challenges unchooseable & $r$ squeezed \emph{after} $(c_{0},c_{\infty})$ absorbed & see next\\
One sumcheck round & false claim $\Rightarrow$ $h$ not identical; Schwartz--Zippel & $\le 3/2^{124}$\\
Opening is committed & path recomputes the root, \texttt{OP\_EQUALVERIFY} & collision\\
Folded row is the committed row & \texttt{hash\_row} reproduces \texttt{PaddingFreeSponge} & preimage\\
Index unchooseable & index bits \emph{are} the path direction bits & ---\\
Answers cannot cancel & random linear combination in powers of $\gamma$ & $\le(t{+}1)/2^{124}$\\
Closing identity & all four EF coefficients; weight taken from the stack & $\le 1/2^{124}$\\
\bottomrule
\end{tabular}
\end{center}

\textbf{Not covered, and so not in that table:} threading the accumulated $(\text{weight},
\text{point})$ pairs between rounds. Both ends exist --- $\sigma'$ is computed and the closing
identity is checked --- but the pairs are not yet carried from round to round, so what the script
guarantees is that each step holds individually, not WHIR's full round-by-round soundness. The gap
between \S\ref{sec:eps}'s $\varepsilon$ and this table is exactly that.

\medskip
\noindent\fbox{\begin{minipage}{0.96\linewidth}\small
\textbf{An assumption that must be named.} The query-term errors above are under the capacity bound
(\texttt{SecurityAssumption::CapacityBound}), which rests on the \emph{unproven} half of
Conjecture 4.12. Under the proved unique-decoding setting this implementation costs 7 to 9 times as
much --- measured at \S\ref{sec:levers}, not read off the paper's table, which is a different
instance --- and script bytes scale with the query count.
\end{minipage}}

\subsection*{Why D1 was ``soundness error $=1$''}

Worth isolating, because it is the line between the first row of that table holding and not. Before
the fix $r$ came from the witness, so the prover could see $r$ before choosing
$(c_{0},c_{\infty})$ --- or choose all three together. For any target $c'$ and any $r\notin\{0,1\}$,
\[
  h(r)=(1-r)c_{0}+r(c-c_{0})+r(r-1)c_{\infty}
\]
is an \textbf{affine} map in $c_{\infty}$ whose coefficient $r(r-1)\neq0$ is invertible, hence
\textbf{surjective} onto $\mathrm{EF}$: pick any $c_{0}$ and solve for the unique $c_{\infty}$ that
hits $c'$. Schwartz--Zippel does not apply --- it bounds ``fix a polynomial, then sample a random
point'', and here the point was not random. Composing across rounds, the closing claim is entirely
the prover's, and the chain's soundness error is 1 rather than something small. That is why the fix
had to bind $r$ to the transcript rather than choose a better constant.

\vspace{1.5em}
\noindent{\color{rule}\rule{\linewidth}{0.6pt}}
\par\vspace{0.5em}
\noindent\footnotesize\color{muted}
Sources: G.~Arnon, A.~Chiesa, G.~Fenzi, E.~Yogev, \emph{STIR: Reed--Solomon Proximity Testing with
Fewer Queries}, ePrint 2024/390 (CRYPTO~'24); and \emph{WHIR: Reed--Solomon Proximity Testing with
Super-Fast Verification}, ePrint 2024/1586. Implementation figures are from the test suite:
\texttt{verifier\_cost}, \texttt{end\_to\_end}, \texttt{verifier},
\texttt{constraint} and \texttt{budget}. The 2162 permutations, 2.9\% transcript share, 8\,967\,342-byte
composed sumcheck chain, per-constraint closing-check sizes and 98 passing tests are all measured on
this machine. What a verifier would cost on-chain, and the
chunking a disprove needs, are the repository README's subject rather than this document's.

\end{document}
